%======================================================================
% 1. 数学环境与格式 (Mathematical Environments & Formatting)
%======================================================================
\usepackage{bm}
\usepackage{amsmath}
\usepackage{amssymb}
\usepackage{mathtools}

%======================================================================
% 2. 基本数学符号 (Basic Mathematical Symbols)
%======================================================================

% -- 微分算子 (Differential Operators) --
\newcommand{\dd}[1]{\mathrm{d} #1}

% -- 导数 (Derivative) --
\newcommand{\ode}[3][1]{%
    \ifnum#1=1
        \frac{\dd{#2}}{\dd{#3}}%
    \else
        \frac{\mathrm{d}^{#1} #2}{\mathrm{d} #3^{#1}}%
    \fi
}

% -- 偏导数 (Partial Derivative) --
\newcommand{\pde}[3][1]{%
    \ifnum#1=1
        \frac{\partial #2}{\partial #3}%
    \else
        \frac{\partial^{#1} #2}{\partial #3^{#1}}%
    \fi
}

% -- 混合偏导数 (Mixed Partial Derivative) --
\newcommand{\mpde}[3]{\frac{\partial^2 #1}{\partial #2 \, \partial #3}}

% -- 向量与算子 (Vectors & Operators) --
\newcommand{\norm}[1]{\lVert #1 \rVert}
\newcommand{\abs}[1]{\lvert #1 \rvert}
\newcommand{\ivec}[1]{\bm{#1}}
\newcommand{\pos}{\ivec{x}}
\newcommand{\vel}{\ivec{u}}
\newcommand{\kvec}{\hat{\ivec{k}}}

% -- 复数与波动 (Complex Numbers & Waves) --
\newcommand{\ee}{\mathrm{e}}
\newcommand{\ii}{\mathrm{i}}
\newcommand{\cc}{\mathrm{c.c.}}
\newcommand{\realpart}{\mathrm{Re}}
\newcommand{\imagpart}{\mathrm{Im}}
\newcommand{\conj}[1]{\overline{#1}}
\newcommand{\complexamp}[1]{\widehat{#1}}
\newcommand{\complexexp}[1]{\ee^{\ii #1}}

%======================================================================
% 3. 物理量定义 (Physical Variables)
%======================================================================

% -- 坐标与时间 --
\newcommand{\timevar}{t}

% -- 地球旋转 --
\newcommand{\earthradius}{a}
\newcommand{\latitude}{\varphi}
\newcommand{\rotationvec}{\ivec{\Omega}}
\newcommand{\coriolis}{f}
\newcommand{\rossby}{\beta}
\newcommand{\ro}{\varepsilon}

% -- 热力学变量 --
\newcommand{\pressure}{p}
\newcommand{\density}{\rho}
\newcommand{\temperature}{T}
\newcommand{\ptemp}{\theta}
\newcommand{\nonadheat}{Q}

% -- 物理常数与参数 --
\newcommand{\gravityconst}{g}
\newcommand{\gasconst}{R_d}
\newcommand{\specheat}{c_p}
\newcommand{\vspecheat}{c_v}
\newcommand{\pref}{p_{\mathrm{ref}}}

% -- 浅水系统变量 --
\newcommand{\freeheight}{\eta}
\newcommand{\bottomheight}{h_b}
\newcommand{\depth}{H}
\newcommand{\gravpotential}{\phi}

% -- 涡度相关 --
\newcommand{\relvort}{\zeta}
\newcommand{\absvort}{\zeta_{\mathrm{abs}}}
\newcommand{\pv}{q}

% -- 尺度分析 --
\newcommand{\hvel}{U}
\newcommand{\vvel}{W}
\newcommand{\hscale}{L}
\newcommand{\vscale}{D}
\newcommand{\timescale}{\tau}

% -- 波动相关 --
\newcommand{\wavefunc}{\psi}
\newcommand{\frequency}{\omega}
\newcommand{\wavenum}{k}

% -- 层结流体相关 --
\newcommand{\bvfreq}{N}
\newcommand{\buoyancy}{B}
\newcommand{\lapse}{\Gamma}
\newcommand{\drylapse}{\Gamma_d}
\newcommand{\scaleheight}{H_s}

% -- 物质导数 --
\newcommand{\matderiv}[1]{\frac{\mathrm{D} #1}{\mathrm{D} \timevar}}

%======================================================================
% 4. 赤道波动特有符号 (Equatorial Wave Specific Symbols)
%======================================================================

% -- 等效深度与特征尺度 --
\newcommand{\eqdepth}{h}
\newcommand{\Leq}{L}
\newcommand{\Teq}{T}
\newcommand{\ceq}{c}

% -- 无量纲变量标记 --
\newcommand{\ndim}[1]{\tilde{#1}}

% -- 组合变量 --
\newcommand{\qvar}{q}
\newcommand{\rvar}{r}

% -- 梯算子 (Ladder Operators) --
\newcommand{\aop}{\hat{a}}
\newcommand{\adop}{\hat{a}^{\dagger}}
\newcommand{\Hop}{\hat{\mathcal{H}}}
\newcommand{\Nop}{\hat{N}}
\newcommand{\Vop}{\hat{V}}
\newcommand{\comm}[2]{\left[ #1, #2 \right]}

% -- 垂直结构算子与函数 --
\newcommand{\vertop}{\hat{\mathcal{L}}}
\newcommand{\horzop}{\hat{\mathcal{M}}}
\newcommand{\vertfunc}{Z}
\newcommand{\vertmode}{m}

% -- 厄米函数与多项式 --
\newcommand{\Hpoly}[1]{H_{#1}}
\newcommand{\hfunc}[1]{\psi_{#1}}
\newcommand{\modenum}{n}

% -- 投影算子 --
\newcommand{\proj}[1]{\mathcal{P}_{#1}}

% -- 湿度参数 --
\newcommand{\moistparam}{\alpha}

%======================================================================
% 5. 特殊波动类型标记 (Wave Type Labels)
%======================================================================

\newcommand{\Kelvin}{\mathrm{Kelvin}}
\newcommand{\Rossby}{\mathrm{Rossby}}
\newcommand{\MRG}{\mathrm{MRG}}
\newcommand{\IGW}{\mathrm{IG}}
\newcommand{\EIG}{\mathrm{EIG}}
\newcommand{\WIG}{\mathrm{WIG}}

%======================================================================
% 6. 常用缩写与标记 (Common Abbreviations & Labels)
%======================================================================

\newcommand{\order}[1]{\mathcal{O}\left(#1\right)}
\newcommand{\const}{\mathrm{const}}
\newcommand{\etal}{\textit{et al.}}
\newcommand{\ie}{\textit{i.e.}}
\newcommand{\eg}{\textit{e.g.}}
\newcommand{\cf}{\textit{cf.}}
\newcommand{\aka}{\textit{a.k.a.}}
